\documentclass[12pt,a4paper,psamsfonts]{amsart}
\usepackage{amssymb,mathtools}
\usepackage{calc}
\setlength{\topmargin}{0cm}
\setlength{\oddsidemargin}{0cm}
\setlength{\evensidemargin}{0cm}
\setlength{\marginparwidth}{0cm}
\setlength{\marginparsep}{0cm}
\setlength{\textheight}{\paperheight - 2in - 35pt}
\setlength{\textwidth}{\paperwidth - 2in}
\setlength{\headheight}{12.5pt}
\setlength{\headsep}{25pt}
\setlength{\footskip}{30pt}
\renewcommand{\baselinestretch}{1.3}
\pagestyle{headings}
\raggedbottom
\setlength{\emergencystretch}{2em}
\usepackage[colorlinks=true,urlcolor=blue,linkcolor=blue,citecolor=blue]{hyperref}
\newtheorem*{theorem}{Theorem}
\theoremstyle{definition}
\newtheorem*{setup}{Setup}
\theoremstyle{plain}
\newtheorem{lemma}{Lemma}[section]
\newtheorem{proposition}[lemma]{Proposition}
\newtheorem{corollary}[lemma]{Corollary}
\newcommand{\Q}{\mathbb Q}
\newcommand{\C}{\mathbb C}
\newcommand{\OO}{\mathcal O}
\newcommand{\mult}{\operatorname{mult}}
\newcommand{\nef}{\operatorname{n}}

\usepackage{tikz-cd}
\newcommand{\R}{\mathbb R}
\newcommand{\Z}{\mathbb Z}
\newcommand{\PP}{\mathbb P}
\DeclareMathOperator{\Pic}{Pic}
\DeclareMathOperator{\Supp}{Supp}
\DeclareMathOperator{\rk}{rk}
\DeclareMathOperator{\ord}{ord}
\DeclareMathOperator{\vol}{vol}
\DeclareMathOperator{\Sym}{Sym}
\DeclareMathOperator{\Bir}{Bir}
\DeclareMathOperator{\codim}{codim}
\def\UrlBreaks{\do\/\do-\do.}
\title{Nef Adjoint Divisors: Framework -- General Dimension}
\date{Temporary draft -- October 11, 2026}
\begin{document}
\maketitle


\section{Theorem and proof}

Fix an integer $d\ge2$. We use the following source inputs,
denoted collectively by $(\mathsf I_d)$.
\begin{itemize}
\setlength{\itemsep}{2pt}
\setlength{\parsep}{0pt}
\item For every projective klt $\Q$-pair $(Y,B)$ of dimension
$j<d$ and every nef $\mathbb{Q}$-Cartier divisor $M$, if $K_Y+B$ is
pseudo-effective and, for some $a\in\Q_{>0}$,
\[
 (K_Y+B+aM)\cdot C>0
\]
for every integral curve through a sufficiently very general point,
then $K_Y+B+cM$ is big for every $c\in\Q_{>0}$.
\item Every projective log canonical $\Q$-pair of dimension at most
$d$ with nef adjoint has semiample adjoint.
\item Every projective $\Q$-factorial generalized klt $\Q$-pair of
dimension at most $d$, with fixed nef $\mathbb{Q}$-Cartier b-divisor
data and pseudo-effective adjoint, admits a terminating choice of
minimal-model program with a nef endpoint. The program extracts no
divisors, preserves $\Q$-factorial generalized klt models, and on a
common resolution the initial adjoint equals the final adjoint plus
an effective divisor exceptional over the final model.
\end{itemize}
The first input is $(P_j)$ of \cite[Section~3.1]{NS}. The last two
are the companion inputs used in \cite[Section~2]{NS}.

\begin{theorem}
Assume $(\mathsf I_d)$. Let $X$ be a projective klt variety over
$\C$ of dimension $d$, with $K_X$ pseudo-effective. Let $D$ be a
nef $\Q$-Cartier divisor such that
\[
 L:=K_X+D\quad\text{is nef},\qquad \nef(X,L)=d.
\]
Then $L$ is big.
\end{theorem}

\begin{proof}
Suppose that $L$ is not big. Since $L$ is nef, $L^d=0$.
After a small projective $\Q$-factorialization, we may assume
that $X$ is $\Q$-factorial. The pullbacks of $K_X,D,L$ retain
the stated properties. Every integral curve through a very
general point of $X$ has positive $L$-degree.

\medskip
\noindent\emph{Adjoint reduction.}
The minimal-model and abundance inputs give $\kappa(K_X)\ge0$.
If $\kappa(K_X)>0$, the positive-Iitaka-dimension reduction in
\cite[Section~3.3]{NS}, using $(P_j)$ for $j<d$, makes $L$ big,
a contradiction. Thus $\kappa(K_X)=0$.

The two-nef-model lemma of \cite[Section~3.2]{NS} gives a
birational contraction $X\dashrightarrow Y$ and $s_0\in\Q_{>0}$
such that $Y$ is projective, $\Q$-factorial and klt,
\[
 K_Y\sim_{\Q}0,\qquad
 N\text{ is nef},\qquad \nef(Y,N)=d,
\]
where $N$ is the transform of $D$. On a common smooth resolution
$X\xleftarrow{p}V\xrightarrow{q}Y$, there is an effective
$q$-exceptional $\Q$-divisor $F$ such that
\[
 p^*(K_X+s_0D)=q^*(K_Y+s_0N)+F.
\]
The divisor $N$ is not big. Indeed, if it were big, this equality
would make $K_X+s_0D$ big. For a rational number
$0<\lambda<\min\{1,1/s_0\}$, the identity
\[
 L=\lambda(K_X+s_0D)+(1-\lambda)K_X+(1-\lambda s_0)D
\]
would then make $L$ big, since the last two summands are
pseudo-effective.

Take the global index-one cover $\nu:Y^{\mathrm{ind}}\to Y$ and
a crepant $\Q$-factorial terminalization
$h:W\to Y^{\mathrm{ind}}$, as in \cite[Section~7]{NS}. Choose an
integer $b>0$ such that $M:=b(\nu\circ h)^*N$ is Cartier. Then
\[
 \begin{gathered}
 W\text{ is projective, $\Q$-factorial and terminal},\qquad
 K_W\sim0,\\
 M\text{ is nef and Cartier},\qquad
 \nef(W,M)=d,\qquad M^d=0.
 \end{gathered}
\]
Full nef dimension is preserved on the common very general
locus. Bigness is detected by the generically finite pullback
$\nu\circ h$, so $M$ is not big. It remains to exclude this
terminal counterexample.

\begin{setup}
\leavevmode
\begin{itemize}
\setlength{\itemsep}{2pt}
\setlength{\parsep}{0pt}
\item Fix an ample Cartier divisor $A$ on $W$.
\item Fix $\epsilon\in\Q_{>0}$ sufficiently small. For $n\ge1$, put
\[
 a_n=((M+A/n)^d)^{1/d},\qquad
 r_n=\max\{1,\lfloor\epsilon/a_n+1/2\rfloor\},\qquad
 P_n=r_n(M+A/n).
\]
Then $P_n$ is ample and $\mathbb{Q}$-Cartier, $r_n\to\infty$, and
$P_n^d\to\epsilon^d$, since $M^d=0$. The index $n$ is independent
of the dimension $d$.
\item Let $p_i:W\times W\to W$ be the projections. Put
\[
 Z=\PP_{W\times W}
       \bigl(p_1^*\OO_W(M)\oplus p_2^*\OO_W(M)\bigr),
 \qquad \pi:Z\to W\times W.
\]
Let $S_1,S_2$ be the two summand sections, and set
\[
 U=\pi^{-1}(W_{\mathrm{reg}}\times W_{\mathrm{reg}})
       \setminus(S_1\cup S_2).
\]
Thus $U$ is smooth and $\dim Z=2d+1$.
\item Set $\xi=c_1(\OO_Z(1))$. For an integer $s>0$, put
\[
 R_{n,s}=(p_1\circ\pi)^*P_n+(p_2\circ\pi)^*P_n+s\xi.
\]
We use the convention
$\pi_*\OO_Z(a)=\Sym^a(p_1^*\OO_W(M)\oplus p_2^*\OO_W(M))$
for $a\ge0$. Whenever $|kR_{n,s}|$ is used, $k>0$ is an integer
with $kP_n$ Cartier.
\item Choose $\epsilon$ sufficiently small that
\[
 (14\epsilon)^d<\tfrac14,\qquad80\epsilon<\tfrac34.
\]
The integer $s$ is chosen after $\epsilon$ and before $n$ tends
to infinity. The multiple $k$ is chosen after $n$.
\end{itemize}
\end{setup}

\medskip
\noindent\textbf{Jet separation.}
Let $V$ be a variety, let $B$ be a Cartier divisor, and let
$z\in V_{\mathrm{reg}}$. For $j\ge0$, $|B|$ separates $j$-jets
at $z$ if
\[
 \forall u\in\OO_{V,z}(B)\quad
 \exists\sigma\in H^0(V,\OO_V(B))\quad
 \sigma_z-u\in\mathfrak m_z^{j+1}\OO_{V,z}(B).
\]

For the terminal counterexample above, we use the following two
lemmas.

\begin{lemma}\label{lem:general-jets}
Assume $(\mathsf I_d)$. Let $W,M$ and $R_{n,s}$ be as in Setup.
There exists an integer $s>0$ such that, for every $n\gg0$,
$|kR_{n,s}|$ separates $(2d+2)k$-jets at a very general point
$z\in U$, for all sufficiently divisible positive integers $k$.
\end{lemma}

This is the large-Seshadri-constant construction of
\cite[Section~5]{NS}, followed by the jet conversion in
\cite[Section~6.7]{NS}.

\begin{lemma}\label{lem:general-nojets}
Assume $(\mathsf I_d)$. Let $W,M$ and $R_{n,s}$ be as in Setup.
For every integer $s>0$ and every $n\gg0$, $|kR_{n,s}|$ does not
separate $(2d+2)k$-jets at a very general point $z\in U$, for any
positive integer $k$ with $kP_n$ Cartier.
\end{lemma}

The terminal application in \cite[Section~6.7]{NS} proves this
obstruction on a smooth resolution $w:\overline W\to W$.
Here $K_{\overline W}\sim E$ for an effective $w$-exceptional
divisor $E$, and an ample perturbation of $w^*P_n$ supplies the
complete-intersection flag. The jet system transfers on the
isomorphism locus. The Frobenius determinant has rank
$p^{2d}$ and its vanishing order satisfies the incompatible bounds
\[
 \frac34p^{2d}\le\ord(\sigma)
 \le\bigl(80\epsilon+O(p^{-1})\bigr)p^{2d}
\]
for arbitrarily large residual characteristics $p$. All complex
data are fixed before reduction.

Fix $s$ supplied by Lemma~\ref{lem:general-jets}, and choose
$n$ sufficiently large for both lemmas. For a sufficiently
divisible $k$, Lemma~\ref{lem:general-jets} gives separation at a
very general point of $U$, whereas
Lemma~\ref{lem:general-nojets} forbids it. This excludes the
terminal counterexample and completes the adjoint reduction.
Hence $L$ is big.
\end{proof}

\medskip
\noindent\textit{Remark.}
The original variety $X$ need not be smooth or have trivial
canonical divisor; no irregularity assumption is imposed.
The condition $K_X$ pseudo-effective is retained from the source.
The trivial-canonical condition occurs only on the auxiliary
terminal model $W$. The two jet lemmas are unchanged in mechanism;
the adjoint reduction precedes their application. No hypothesis
$\kappa(D)=-\infty$ is needed. The full-nef-dimension target is
bigness of $K_X+D$; this does not assert nonvanishing for an
arbitrary nef adjoint of smaller nef dimension.

\medskip
\noindent\textbf{Draft status.}
This is a conditional framework following \cite[Sections~3--7]{NS}.
The source inputs $(\mathsf I_d)$ and the inherited geometric and
Frobenius arguments have not been independently verified here.
The two jet lemmas record the source proof obligations; they are
not independent proofs supplied by this framework.

\begin{thebibliography}{99}
\bibitem{NS}
\emph{Numerical Semiampleness of Nef Adjoint Divisors},
OpenAI mathematics collection, October 3, 2026,
Sections~3--7.
\href{https://github.com/openai/math/blob/main/preprints/Numerical-Semiampleness-of-Nef-Adjoint-Divisors-October-3-2026/paper.pdf}{Source manuscript}.
\end{thebibliography}

\end{document}
